\section{Analisis Masalah}

Sebelum artefak dirancang, persoalan penambangan kebijakan ABAC pada \textit{edge-IoT} diuraikan menjadi sejumlah \textit{design tension} yang harus dijawab.
Analisis ini menjadi dasar penurunan kebutuhan dan keputusan rancangan dari algoritma CoDA-EDA yang diusulkan.

\begin{enumerate}
    \item Keterbatasan sumber daya perangkat \textit{edge} -- Perangkat \textit{edge-IoT} beroperasi dengan memori dan daya komputasi yang terbatas. Algoritma metaheuristik berpopulasi eksplisit, seperti BOA, ACO, dan UBK, memelihara sekumpulan kandidat solusi yang jejak memorinya tumbuh seiring ukuran populasi, sehingga tidak layak dijalankan langsung pada perangkat semacam itu. Dari masalah ini muncullah persoalan bagaimana mencari kebijakan berkualitas tanpa memelihara populasi solusi yang membengkak.
    \item Ketegangan antara kesadaran dependensi dan hemat memori -- Atribut ABAC di dunia nyata itu saling berkorelasi (mis. peran-unit). Keluarga \textit{compact} univariat hemat memori tetapi mengasumsikan independensi penuh antar-atribut, sehingga kualitasnya menurun pada ruang berdimensi tinggi. Sebaliknya, algoritma yang memodelkan dependensi penuh (jaringan bayes pada BOA/iBOA) memerlukan representasi struktur yang mahal (hingga $O(D^2)$). Dari masalah ini muncullah persoalan bagaimana menangkap keterkaitan antar-atribut dengan biaya memori yang tetap hemat.
    \item Dinamika kebijakan dan biaya pembaruan -- Lingkungan IoT bersifat dinamis sehingga kebijakan perlu diperbarui berkala mengikuti perubahan distribusi data (\textit{drift}). Menambang ulang dari nol, atau membangun ulang struktur dependensi pada setiap siklus sebagaimana pendekatan konvensional akan memakan waktu dan energi yang tidak efisien pada perangkat \textit{edge}, sehingga muncul suatu persoalan bagaimana memperbarui kebijakan secara berkala dengan biaya komputasi minimal tanpa menurunkan kualitas.
    \item Kompromi kualitas yang proporsional -- Pada konteks \textit{edge}, sedikit kompromi kualitas penambangan statis dapat diterima asalkan memori stabil dan pembaruan dilakukan secara efisien. Namun, kualitas kebijakan yang dihasilkan harus tetap kompetitif terhadap algoritma sekelas agar artefak tetap bermanfaat. Dari situlah muncul persoalan yakni bagaimana menjaga kualitas tetap setara sambil memenuhi ketiga kendala di atas. 
\end{enumerate}

Keempat persoalan di atas dapat ditelusuri secara konkret pada enam algoritma pembanding yang telah diidentifikasi, yang mana masing-masing memenuhi sebagian kebutuhan tetapi gagal pada sebagian lain, sehingga tidak ada satu pun yang sekaligus ringan, sadar-dependensi, dan inkremental untuk \textit{edge}.
Pertama, pada keluarga algoritma \textit{compact}, seperti \textit{compact-GA} dan \textit{compact-PSO}.
Keduanya memenuhi kebutuhan memori (jejak O(D) yang datar terhadap populasi), tetapi mengasumsikan independensi penuh antar-atribut.
Pada ruang kebijakan berdimensi tinggi yang kaya dengan keterkaitan antar-atribut, asumsi ini dapat menurunkan kualitas karena kombinasi atribut yang selaras tidak terwakili.
Masalah utama yang ada pada keluarga algoritma \textit{compact} ini selaras dengan peroalan kedua, yakni buta terhadap dependensi.

Yang kedua, pada kelompok algoritma EDA yang sadar dependensi, seperti BOA dan iBOA.
Keduanya sadar dependensi, tetapi bermasalah pada sisi yang berbeda.
BOA memelihara populasi eksplisit dan membangun ulang jaringan bayes pada setiap generasi, sehingga jejak memorinya tumbuh dan biaya pembaruannya menjadi tinggi, sehingga menjadikannya tidak layak untuk perangkat \textit{edge}.
iBOA memperbaiki sisi memori (jejaknya datar melalui statistik inkremental) dan tetap sadar dependensi, sehingga menjadi pembanding terdekat dengan algoritma yang diusulkan.
Namun ia memperbarui struktur secara terus-menerus pada setiap langkah pencarian, tidak dipicu oleh perubahan data atau \textit{drift}, dan tanpa mekanisme \textit{warm-start}.
Akibatnya, pada pembaruan berkala, iBOA melakukan pembelajaran ulang struktur yang tidak perlu ketika data sebenarnya stabil.

Yang ketiga adalah keluarga algoritma metaheuristik berpopulasi eksplisit, seperti ACO dan UBK.
Keduanya mencapai kualitas tinggi pada penambangan statis, tetapi berbasis populasi eksplisit/feromon yang jejak memorinya membengkak seiring ukuran populasi, dan dirancang untuk eksekusi \textit{batch/offline} pada sumber daya yang memadai, bukan pembaruan inkremental di \textit{edge}.

Dari ketiga keluarga algoritma tersebut, bisa disimpulkan bahwa tidak ada satu pun algoritma pembanding yang memenuhi ketiga kebutuhan inti secara serentak.
iBOA sebenarnya adalah algoritma yang paling mendekati, tetapi gagal pada pembaruan berkala yang efisien.
Kesenjangan inilah yang secara langsung memotivasi rancangan CoDA-EDA.
Ringkasannya disajikan pada Tabel \ref{tab:analisis-masalah}

\begin{table}[htbp]
\centering
\footnotesize
\setlength{\tabcolsep}{4pt}
\renewcommand{\arraystretch}{1.15}
\caption{Analisis Masalah pada Algoritma Pembanding terhadap Kebutuhan Inti Edge-IoT}
\label{tab:analisis-masalah}
\begin{tabular}{@{} 
  >{\raggedright\arraybackslash}p{3.2cm} 
  >{\centering\arraybackslash}p{1.4cm} 
  >{\centering\arraybackslash}p{1.8cm} 
  >{\centering\arraybackslash}p{1.8cm} 
  >{\raggedright\arraybackslash}p{5.5cm} 
@{}}
\toprule
Algoritma (keluarga) & Sadar dependensi & Memori datar thd populasi & Pembaruan inkremental efisien & Masalah utama di edge-IoT \\
\midrule
compact-GA (compact) & tidak & ya & tidak & Buta dependensi $\rightarrow$ kualitas turun pada dimensi tinggi \\
compact-PSO (compact) & tidak & ya & tidak & Independensi penuh antar-atribut, sama seperti compact-GA \\
BOA (EDA penuh) & ya & tidak & tidak & Populasi eksplisit + bangun ulang struktur tiap generasi $\rightarrow$ boros \\
iBOA (EDA inkremental) & ya & ya & tidak & Relearn struktur kontinu, tak sadar-drift, tanpa warm-start \\
ACO (berpopulasi) & ya\textsuperscript{*} & tidak & tidak & Populasi/feromon tumbuh; batch/offline, bukan edge \\
UBK (berpopulasi) & ya\textsuperscript{*} & tidak & tidak & Berpopulasi eksplisit $\rightarrow$ memori tumbuh; tidak inkremental \\
CoDA-EDA (usulan) & ya (orde-1) & ya & ya & memenuhi ketiga kebutuhan inti \\
\bottomrule
\end{tabular}
\par\vspace{0.5em}
{\footnotesize ya\textsuperscript{*} = keterkaitan atribut ditangani secara implisit melalui jejak feromon/konstruksi solusi, bukan model dependensi eksplisit.}
\end{table}